\section{Mở đầu}

\subsection{Bối cảnh và động lực}
Nhận dạng mẫu là một lĩnh vực nền tảng trong thị giác máy tính và trí tuệ nhân tạo, tập trung vào việc xây dựng các phương pháp giúp máy tính phát hiện, biểu diễn, phân loại hoặc xác minh đối tượng dựa trên dữ liệu quan sát được. Trong đồ án này, nhóm lựa chọn một chủ đề nhận dạng cụ thể để nghiên cứu theo hai hướng song song: nền tảng lý thuyết từ giáo trình và hướng tiếp cận hiện đại từ một bài báo khoa học tiên tiến.

\subsection{Mục tiêu đồ án}
Báo cáo hướng đến các mục tiêu chính sau:
\begin{itemize}
    \item Trình bày có hệ thống kỹ thuật hoặc phương pháp cốt lõi trong chương sách đã chọn.
    \item Phân tích một bài báo SOTA có liên quan, bao gồm động lực, ý tưởng chính, phương pháp, thực nghiệm và thảo luận.
    \item Làm rõ mối liên hệ giữa kiến thức nền tảng trong giáo trình và hướng nghiên cứu hiện đại.
    \item Triển khai demo hoặc thực nghiệm minh họa, phân tích mã nguồn, thiết lập thực nghiệm và kết quả thu được.
    \item Đánh giá ưu điểm, hạn chế và đề xuất hướng mở rộng cho phương pháp được khảo sát.
\end{itemize}

\subsection{Phạm vi và đối tượng nghiên cứu}
\textbf{Chủ đề nhận dạng được chọn:} \textit{[Điền chủ đề cụ thể, ví dụ: nhận dạng khuôn mặt, nhận dạng vân tay, nhận dạng dáng đi, nhận dạng mống mắt, v.v.]}.

\textbf{Chương sách giáo trình:} \textit{[Điền tên chương, tên sách: Handbook of Biometrics hoặc Handbook of Face Recognition 2nd Edition]}.

\textbf{Bài báo SOTA:} \textit{[Điền tên bài báo, hội nghị/tạp chí, năm công bố và liên kết mã nguồn nếu có]}.

\subsection{Cấu trúc báo cáo}
Phần còn lại của báo cáo được tổ chức như sau. Mục 2 trình bày kiến thức nền tảng cần thiết. Mục 3 tóm tắt và phân tích chương sách giáo trình. Mục 4 phân tích bài báo khoa học tiên tiến và so sánh với nền tảng lý thuyết. Mục 5 trình bày phần thực hành, bao gồm demo, phân tích mã nguồn, thiết lập thực nghiệm và kết quả. Mục 6 đưa ra đánh giá, nhận định, phân công nhiệm vụ và kết luận.
